% Generated draft: requires mathematical and source review.
\documentclass[11pt]{article}
\usepackage[margin=25mm]{geometry}
\usepackage{fontspec}
\setmainfont{Latin Modern Roman}
\usepackage{amsmath,amssymb,mathtools,unicode-math}
\setmathfont{Latin Modern Math}
\usepackage{enumitem,needspace,xurl,hyperref}
\setlength{\parindent}{0pt}
\setlength{\parskip}{5pt}
\setlength{\emergencystretch}{3em}
\newcommand{\N}{\mathbb{N}}
\newcommand{\Z}{\mathbb{Z}}
\newcommand{\Q}{\mathbb{Q}}
\newcommand{\R}{\mathbb{R}}
\newcommand{\C}{\mathbb{C}}
\newcommand{\F}{\mathbb{F}}
\newcommand{\A}{\mathbb{A}}
\newcommand{\PP}{\mathbb{P}}
\newcommand{\E}{\mathbb{E}}
\newcommand{\eps}{\varepsilon}
\newcommand{\dd}{\,\mathrm{d}}
\begin{document}
\section*{UnsolvedMath ID: 700784}
\textbf{Source problem code:} catalogue entry
\subsection{Definitions and mathematical statement}
All groups considered are finite, and classes of groups are understood to be closed under isomorphism. A formation is a nonempty class \(\mathcal F\) closed under homomorphic images and finite subdirect products: if \(G\leq G_1\times\cdots\times G_n\), with \(n\geq1\) and \(G_i\in\mathcal F\), and every coordinate projection of \(G\) is surjective, then \(G\in\mathcal F\). A formation is hereditary if every subgroup of each of its members also belongs to it. A group \(G\) is soluble if its derived series, defined by \(G^{(0)}=G\) and \(G^{(i+1)}=[G^{(i)},G^{(i)}]\), reaches the trivial group after finitely many steps. A formation is soluble if all its members are soluble.

For a finite group \(X\), let
\[
\operatorname{form}(X)=\bigcap\{\mathcal C:\mathcal C\text{ is a formation and }X\in\mathcal C\}.
\]
Here a formation is called one-generated if it equals \(\operatorname{form}(X)\) for some finite group \(X\).

The classification problem proposed by A. N. Skiba is to determine all soluble hereditary formations \(\mathcal F\) such that
\[
\begin{gathered}
\forall X\in\mathcal F,\qquad \mathcal F\ne\operatorname{form}(X),\\
\forall\mathcal H\subsetneq\mathcal F,\qquad
\bigl(\mathcal H\text{ is a hereditary formation}\bigr)
\Longrightarrow
\bigl(\exists X\in\mathcal H:\mathcal H=\operatorname{form}(X)\bigr).
\end{gathered}
\]
Thus \(\mathcal F\) is not one-generated, but every proper hereditary subformation is one-generated. The trivial formation is included among formations under these definitions.

The supplied excerpt does not define the generating operation. The formulation above uses generation in the class of all formations. If generation within the class of hereditary formations was intended instead, the generating closure in both conditions must be changed accordingly; the supplied record alone does not resolve this distinction.
\subsection{Short English statement}
Classify the hereditary formations of finite soluble groups that cannot be generated by one finite group, although every proper hereditary subformation can. Generation is interpreted above as ordinary formation generation; the supplied excerpt leaves this convention unstated.
\subsection{Sources}
\begin{itemize}
\item Local UnsolvedMath registry, supplied complete source record: permanent ID \(700784\), exact source problem code \texttt{catalogue entry}, title \emph{Kourovka Notebook Problem 11.92}. Upstream revision: \texttt{d371ee917ac9647fbb06cd6fc529deec6f3dda82}.
\item The supplied record attributes the problem to A. N. Skiba, proposed in 1990, and identifies its source as \emph{The Kourovka Notebook (Unsolved Problems in Group Theory)}, 11th issue (1990), Problem 11.92; it also identifies the 21st issue edition (2026), arXiv:1401.0300v48. The collection description names E. I. Khukhro and V. D. Mazurov as editors.
\item The local record labels the problem \texttt{open}. This is a supplied registry classification, not an independently verified mathematical or literature assessment. The cited source and its terminology have not been independently checked for this draft.
\end{itemize}
\end{document}
